Stolové a kartové hry sú dlhodobo súčasťou ľudstva. Slúžia na pobavenie, spríjemnenie dlhých chvíľ alebo aj edukáciu. Taktiež majú sociálny aspekt, ktorý preukázateľne pomáha mentálnemu zdraviu človeka. Mobilné zariadenia, alebo aj smartfóny, sú každodennou súčasťou našich životov a využívame ich na rôzne aktivity. Spojenie stolových hier a smartfónov môže priniesť veľa výhod, ktoré tradičné papierové manuály a pravidlá nedokážu. Prehrávanie tematickej hudby, časové obmedzenia alebo aj rýchle skenovanie herných objektov prinášajú benefity do tohto prostredia.

Táto bakalárska práca sa zaoberá využitím počítačového videnia na mobilných zariadeniach s cieľom uľahčiť hranie stolových hier. Tento trend je na vzostupe vďaka čoraz vyššiemu výkonu mobilných zariadení. Vďaka tomu je možné vytvárať aplikácie, ktoré rozpoznávajú herné objekty, napríklad karty, figúrky, ale aj kocky, a poskytujú hráčom dodatočné informácie počas hrania hier.

Cieľom práce je navrhnúť a implementovať znovupoužiteľný pracovný postup pre vývojárov androidových aplikácií, ktoré využívajú objektovú detekciu. Výsledkom tejto práce je knižnica a desktopová aplikácia, ktoré proces tvorby týchto aplikácií zjednodušujú.

Prvou časťou je androidová knižnica CVLiBG, napísaná v jazyku Kotlin, ktorá zjednodušuje prácu s kamerou a implementovanie detekčných modelov do aplikácií. Knižnica ďalej pomáha zobrazovať výsledky tohto procesu a umožňuje moderný kotlinový spôsob spracovania týchto výsledkov. Knižnica je navrhnutá modulárne, aby vývojári mohli použiť pripravené komponenty alebo ich nahradiť vlastnou implementáciou.

Modely detekovania objektov sú všeobecne rozšírené, no špecifické modely pre konkrétne hry alebo herné objekty nemusia existovať. Táto práca sa pokúša tento problém vyriešiť, hlavne pre menších vývojárov, pomocou druhej časti tejto práce, desktopovej aplikácie Dataset Creator. Aplikácia slúži na tvorbu obrazových datasetov potrebných na trénovanie modelov objektovej detekcie. Aplikácia podporuje anotáciu obrázkov, generovanie syntetických dát, export datasetov a trénovanie modelov pomocou predtrénovaných modelov.

Navrhnuté riešenie bolo overené na demonštračnej aplikácii pre kartovú hru Bang!. Pre túto hru bol vytvorený model na rozpoznávanie herných kariet Bang!. Dataset obrazových dát bol vytvorený z nazbieraných fotografií a označený v aplikácii. Ďalej boli na dataset aplikované syntetické dáta, ktoré mali za úlohu zvýšiť robustnosť modelu. Finálny model bol taktiež natrénovaný v aplikácii. Tento model bol následne integrovaný do mobilnej aplikácie pomocou knižnice CVLiBG, pričom samotná aplikácia sa nemusela zameriavať na proces detekcie, ale len na zbieranie výsledkov. Výsledná aplikácia dokáže pomocou kamery mobilného telefónu rozpoznávať vybrané karty a zobrazovať používateľovi informácie súvisiace s detegovanými objektmi.

Hlavným prínosom práce je ucelený pracovný postup od prípravy dát až po ich použitie v mobilnej aplikácii. CVLiBG umožňuje vývojárom sústrediť sa najmä na používateľské rozhranie a logiku aplikácie, zatiaľ čo technicky náročné časti, ako práca s kamerou a objektová detekcia, sú abstrahované v knižnici. Dataset Creator zároveň zjednodušuje prípravu dát a trénovanie modelov pre konkrétne herné objekty. Výsledky práce ukazujú, že mobilná objektová detekcia môže byť prakticky využiteľná pri tvorbe aplikácií podporujúcich hranie stolových hier, hoci jej úspešnosť závisí od kvality datasetu, výkonu zariadenia a zložitosti snímanej scény.